\documentclass[11pt]{article}
\usepackage[T1]{fontenc}
\usepackage[utf8]{inputenc}
\usepackage{newpxtext}
\usepackage{amsmath,amsthm}
\usepackage{newpxmath}
\usepackage[a4paper,margin=28mm,headheight=14pt]{geometry}
\usepackage{microtype}
\usepackage{enumitem}
\usepackage{needspace}
\usepackage{xcolor}
\usepackage{xurl}
\usepackage[unicode,bookmarksnumbered]{hyperref}
\usepackage[nameinlink,noabbrev]{cleveref}
\usepackage{fancyhdr}
\definecolor{linkblue}{RGB}{26,66,108}
\hypersetup{colorlinks=true,linkcolor=linkblue,citecolor=linkblue,urlcolor=linkblue,
 pdftitle={Sharp cardinal bounds for cup products on locally profinite spaces},
 pdfsubject={A sharp answer to Aoki's Question 2.14, with a Lean formalization}}
\urlstyle{tt}
\setlength{\emergencystretch}{1.2em}
\setlength{\parindent}{1.25em}
\setlength{\parskip}{2pt}
\setlist[enumerate]{itemsep=4pt,topsep=6pt}
\setlist[itemize]{itemsep=3pt,topsep=5pt}
\pagestyle{fancy}
\fancyhf{}
\fancyhead[L]{\small Sharp cardinal bounds for cup products}
\fancyhead[R]{\small\thepage}
\renewcommand{\headrulewidth}{0.3pt}
\fancypagestyle{plain}{\fancyhf{}\fancyfoot[C]{\thepage}\renewcommand{\headrulewidth}{0pt}}
\theoremstyle{plain}
\newtheorem{theorem}{Theorem}[section]
\newtheorem{proposition}[theorem]{Proposition}
\newtheorem{lemma}[theorem]{Lemma}
\newtheorem{corollary}[theorem]{Corollary}
\theoremstyle{definition}
\newtheorem{definition}[theorem]{Definition}
\newtheorem{question}[theorem]{Question}
\theoremstyle{remark}
\newtheorem{remark}[theorem]{Remark}
\newtheorem{example}[theorem]{Example}
\BeforeBeginEnvironment{theorem}{\Needspace{17\baselineskip}}
\BeforeBeginEnvironment{lemma}{\Needspace{8\baselineskip}}
\BeforeBeginEnvironment{proposition}{\Needspace{8\baselineskip}}
\BeforeBeginEnvironment{corollary}{\Needspace{7\baselineskip}}
\newcommand{\F}{\mathbf F_2}
\newcommand{\Z}{\mathbf Z}
\newcommand{\Ab}{\mathbf{Ab}}
\DeclareMathOperator{\Hom}{Hom}
\DeclareMathOperator{\Ext}{Ext}
\DeclareMathOperator{\Sh}{Sh}
\DeclareMathOperator{\pd}{pd}
\DeclareMathOperator{\supp}{supp}
\newcommand{\leanmark}[1]{\unskip\textsuperscript{\hyperref[lean:#1]{\textsf{[L\ref*{lean:#1}]}}}}
\newcounter{leanentry}
\newcommand{\leanentry}[2]{\par\addvspace{9pt}\Needspace{5\baselineskip}\refstepcounter{leanentry}\label{lean:#1}\noindent\textbf{L\theleanentry. #2}\par\nobreak}
\newenvironment{leansources}{\begingroup\small\raggedright\begin{itemize}[leftmargin=1.3em,itemsep=2pt,topsep=3pt]}{\end{itemize}\endgroup}
\title{\bfseries Sharp cardinal bounds for cup products\\on locally profinite spaces}
\author{}
\date{}
\begin{document}
\thispagestyle{plain}
\begin{center}
{\Large\bfseries Sharp cardinal bounds for cup products\\[3pt]
on locally profinite spaces\par}
\end{center}
\vspace{0.5em}
\begin{abstract}
For every positive integer \(r\), we construct a locally profinite Hausdorff
space of cardinality and weight \(\aleph_r\) with a degree-one class over
\(\mathbf F_2\) whose \(r\)-th cup power is nonzero. These sizes are minimal:
nonzero degree-\(r\) sheaf cohomology, even with arbitrary abelian sheaf
coefficients, forces both cardinality and weight at least \(\aleph_r\).
This answers Question 2.14 of Aoki. The construction uses a punctured
involution quotient. Its cup-power nonvanishing is reduced to an elementary
finite-difference argument: continuity gives finite exceptional sets,
cardinal growth lets us avoid them uniformly, and symmetry makes induction
possible. A separate argument with compact-open covers proves the universal
lower bound. We give the full proof and a declaration-level guide to its
Lean formalization.
\end{abstract}
\section{The question and the answer}
\label{sec:introduction}

A locally profinite space has a basis of compact-open sets. Locally, such
sets are particularly well behaved for sheaf cohomology: sections on a
compact open have no higher derived functors. Nevertheless, assembling
many compact opens can produce nonzero cohomology, and even nonzero
products of positive-degree classes. How large must the space be for
this to happen?

In \cite[Theorem C]{Aoki}, Aoki constructs, for each \(n\geq0\), a locally
profinite space \(U\) of cardinality \(\aleph_{2n+1}\) with classes
\(\eta_0,\ldots,\eta_n\in H^1(U;\F)\) such that

\[
\eta_0\smile\cdots\smile\eta_n\neq0.
\]

He then asks the following question.

\begin{quote}
\textbf{Aoki's Question 2.14.}
Is Theorem C optimal in terms of the weight (or cardinality) of \(U\)?
\end{quote}

We prove that the sharp bound for a product of \(n+1\) degree-one classes
is \(\aleph_{n+1}\), for both cardinality and weight. All factors in the
product can in fact be the same class. To keep the cohomological degree
visible, we write \(r=n+1\) throughout the proof.

A \emph{locally profinite Hausdorff space} will mean a locally compact,
Hausdorff, totally disconnected space. A compact Hausdorff totally
disconnected space is called profinite. Our examples are a profinite
space with one point removed, so they also satisfy the convention used
in \cite{Aoki}. The weight \(w(V)\) is the least cardinality of a basis for the topology.
All cohomology is ordinary sheaf cohomology; a constant
coefficient such as \(\F\) means the associated constant sheaf. Powers
of a degree-one class are cup powers. We work with the axiom of choice;
no continuum hypothesis is assumed.

\begin{theorem}\label{thm:main}
For every integer \(r\geq1\), there exist a locally profinite Hausdorff
space \(U_r\) and a class \(\eta_r\in H^1(U_r;\F)\) such that

\[
|U_r|=w(U_r)=\aleph_r,
\qquad \eta_r^r\neq0,
\qquad \eta_r^{r+1}=0.
\]

Moreover, for every abelian sheaf \(\mathcal M\) on \(U_r\),

\[
H^q(U_r;\mathcal M)=0\qquad(q>r).
\]

Conversely, for every locally profinite Hausdorff space \(V\) and every
abelian sheaf \(\mathcal M\) on \(V\),

\[
H^r(V;\mathcal M)\neq0
\quad\Longrightarrow\quad
|V|\geq\aleph_r\quad\text{and}\quad w(V)\geq\aleph_r.
\]

Thus \(\aleph_r\) is both the least cardinality and the least weight of a
locally profinite Hausdorff space supporting a nonzero product of \(r\)
degree-one constant-\(\F\) classes.\leanmark{main}
\end{theorem}

The theorem shows that Aoki's bound is optimal when \(n=0\), and can be
lowered for every \(n\geq1\). It also gives examples of cohomological
dimension exactly \(r\), where dimension is taken over all abelian sheaf
coefficients.

\subsection{Reading the paper and the formalization}

The proof is accompanied by a Lean development. Small references such as
\leanmark{main} point to the source concordance in
\cref{app:lean}, where each mathematical step is matched with explicit
files, declarations, and line numbers. These references can be skipped
on a first reading. The mathematical argument does not require knowledge
of Lean. \Cref{sec:formalization} explains precisely what is checked and
how to reproduce the verification.

\subsection{The proof in five steps}
\label{subsec:strategy}

There are two logically independent tasks: attaining the size
\(\aleph_r\), and excluding every smaller size. The first task contains
the main construction. Its guiding idea is to build a cocycle for which
being a boundary would force an impossible equation between functions.

\begin{enumerate}
\item \textbf{Put symmetry on coordinates of increasing sizes.}
Take \(r+1\) coordinates with underlying sizes
\(\aleph_0,\ldots,\aleph_r\). Each coordinate has a bit in \(\F\)
and a point at infinity. Identify points that differ by flipping every
bit, and remove the all-infinite point. This is \(U_r\). It has a cover
by \(r+1\) opens, according to which coordinate is finite.

\item \textbf{Write the class and its power explicitly.}
On the \(i\)-th open, choose the representative whose \(i\)-th bit is
zero. On an overlap, the difference of the two choices is
\(z_{ij}=s_i+s_j\). These differences give \(\eta_r\), and its top
power is represented by

\[
F_r=\prod_{j=0}^{r-1}(s_j+s_{j+1}).
\]

The intersections in the cover are disjoint unions of compact opens.
This makes their free sheaves projective, so an explicit finite
resolution computes ordinary cohomology and its cup product.

\item \textbf{Extract the restrictions on a possible boundary.}
If \(F_r\) were a boundary, it would be a sum
\(f_0+\cdots+f_r\). Each \(f_i\) would be invariant under the
simultaneous bit flip. It would also be constant outside finitely many
points in its \(i\)-th coordinate, with the other coordinates fixed:
this is exactly what continuity at infinity forces.

\item \textbf{Use the last coordinate to lower the degree.}
There are only \(\aleph_{r-1}\) configurations of the earlier
coordinates. We can therefore choose one base point in the last
coordinate that avoids every finite exceptional set for \(f_r\).
Adding the values at its two bits kills \(f_r\) and changes \(F_r\)
into \(F_{r-1}\). The other summands retain their symmetry and
almost-constancy. Induction reaches a two-coordinate contradiction:
an almost-constant function cannot differ by \(1\) from its bit flip
everywhere.

\item \textbf{Bound the cohomology of every smaller space.}
Independently, a cover by at most \(\aleph_n\) compact opens forces
cohomology to vanish above degree \(n\). Countable covers can be
disjointified; passing from one aleph to the next uses a continuous
transfinite filtration and raises the projective-dimension bound by
at most one. Both cardinality and weight control the size of a
compact-open cover. This proves the lower bounds.
\end{enumerate}

The roles of the three ingredients in the nonvanishing argument are
worth keeping distinct. \emph{Continuity} produces finite exceptional
sets. \emph{Cardinal growth} permits a choice avoiding all of them at
once. \emph{Symmetry} makes the two evaluations fit together into an
induction. The projective resolution connects this elementary argument
to sheaf cohomology.

\subsection{The first example as a guide}

For \(r=1\), the two coordinate sets have sizes \(\aleph_0\) and
\(\aleph_1\), and the cocycle is simply \(s_0+s_1\). A hypothetical
boundary would split it into two invariant functions, the first
almost constant in coordinate zero and the second almost constant in
coordinate one. Countably many finite exceptional sets in the second
coordinate cannot exhaust its uncountable set of base points. Evaluate
at both bits of a base point outside their union. The second function
cancels, and the resulting slice of the first would have to change by \(1\)
under bit flip.
Far enough from its own finite exceptional set, however, it is constant.
This contradiction is the base case for every higher cup power.

We now carry out the construction in \cref{sec:geometry}, the
cohomological calculation in \cref{sec:resolution,sec:cup}, and the
nonvanishing argument in \cref{sec:nonvanishing}. The independent
optimality proof occupies \cref{sec:optimality}.

\section{The space and its compact-open pieces}
\label{sec:geometry}

Fix \(r\geq1\). Choose discrete sets \(A_i\) with
\(|A_i|=\aleph_i\), for \(0\leq i\leq r\), and set

\[
D_i=A_i\times\F,
\qquad K_i=D_i\sqcup\{\infty_i\},
\qquad K=\prod_{i=0}^{r}K_i.
\]

Give \(K_i\) the one-point compactification topology: every point of
\(D_i\) is isolated, and the neighborhoods of \(\infty_i\) contain
all but finitely many points of \(D_i\). The spaces \(K_i\) and their
finite product \(K\) are profinite.

Let \(\tau\) flip the bit in every finite coordinate and fix each
point at infinity:

\[
\tau(a,s)=(a,s+1),\qquad \tau(\infty_i)=\infty_i.
\]

This defines a continuous involution of \(K\). Its only fixed point
is \(p=(\infty_0,\ldots,\infty_r)\), because any finite coordinate
changes its bit. Write

\[
q:K\longrightarrow\overline K=K/\langle\tau\rangle,
\qquad U_r=\overline K\setminus\{q(p)\}.
\]

Taking the compact quotient first makes the topology of the example
particularly transparent. The quotient is compact Hausdorff: its orbit
relation is the union of the diagonal and the graph of a homeomorphism,
hence is closed. It is also zero-dimensional. Indeed, in the inverse
image of any open neighborhood of an orbit, choose a clopen neighborhood
of one representative; its union with its translate is invariant and
clopen, and descends to a clopen neighborhood in the quotient. The
orbit map is open, since the saturation of an open set \(O\) is
\(O\cup\tau(O)\). Thus \(\overline K\) is profinite and \(U_r\)
is locally profinite Hausdorff.\leanmark{construction}

\subsection{Choosing a representative on each chart}

Let \(W_i\subseteq U_r\) consist of the orbits whose \(i\)-th
coordinate is finite. These opens cover \(U_r\): the only orbit with
no finite coordinate was removed. On \(W_i\), each orbit has exactly
one representative whose \(i\)-th bit is zero. This representative
depends continuously on the orbit, because the bit-zero slice is open
and the quotient map restricts to an open continuous bijection from
that slice onto \(W_i\). We call this the \emph{gauge choice} on
\(W_i\).\leanmark{charts}

\begin{lemma}\label{lem:cells}
For every nonempty \(I\subseteq\{0,\ldots,r\}\), the intersection
\(W_I=\bigcap_{i\in I}W_i\) is a disjoint union of compact-open
subsets of \(U_r\).
\end{lemma}

\begin{proof}
Choose an anchor \(a\in I\) and use its gauge choice. Prescribe every
finite coordinate \((a_i,s_i)\in D_i\) for \(i\in I\), with
\(s_a=0\), and leave all coordinates outside \(I\) unrestricted.
This gives a clopen cylinder in \(K\). Its image in the quotient is
compact and open, and avoids \(q(p)\).

These images are disjoint for distinct prescriptions. If two meet,
the corresponding representatives either agree or are exchanged by
\(\tau\). Exchange is impossible because both anchor bits are zero;
agreement forces identical prescriptions. Every orbit in \(W_I\)
has one such prescription, so the images cover \(W_I\).
\end{proof}

These compact-open pieces explain why the cover will compute
cohomology. The intersections need not themselves be compact; it is
their decomposition into disjoint compact pieces that matters.

\subsection{Computing cardinality and weight}

\begin{proposition}\label{prop:size}
The space \(U_r\) satisfies \(|U_r|=w(U_r)=\aleph_r\).
\end{proposition}

\begin{proof}
The finite product \(K\) has cardinality \(\aleph_r\), the size of its
largest factor. A quotient and then a subspace cannot increase
cardinality, so \(|U_r|\leq\aleph_r\).

For the weight estimate, recall that an infinite compact Hausdorff
space \(T\) has \(w(T)\leq|T|\). Here is a useful direct proof.
For each ordered pair of distinct points \(x,y\), choose disjoint
neighborhoods \(O_{xy}\) of \(x\) and \(V_{xy}\) of \(y\).
Given an open neighborhood \(N\) of \(x\), finitely many of the \(V_{xy}\)
cover the compact set \(T\setminus N\). The intersection of the
corresponding \(O_{xy}\) is contained in \(N\). These finite
intersections, for all \(x\), form a basis with at most \(|T|\)
elements.

An open quotient map does not increase weight, because the images of
basis elements form a basis. Subspaces do not increase weight either.
Consequently,

\[
w(U_r)\leq w(\overline K)\leq w(K)\leq\aleph_r.
\]

For both reverse inequalities, consider the axis whose only finite
coordinate is the last one. Every orbit there has a unique
representative with last coordinate \((a,0)\), so the axis is a copy
of \(A_r\). It is a discrete subspace: the open condition that the
last base coordinate equal \(a\) isolates its corresponding axis
point. We obtain \(\aleph_r\) distinct points, and any basis for
\(U_r\) must contain distinct neighborhoods isolating these points
in the axis. Thus both cardinality and weight are at least
\(\aleph_r\).\leanmark{size}
\end{proof}

\section{A finite resolution that computes cohomology}
\label{sec:resolution}

The geometry above serves two purposes. It bounds the size of the
space, and it provides a cover with projective free sheaves on all its
intersections. We spell out the latter point so that the passage from
functions to ordinary sheaf cohomology is explicit.

\subsection{Free sheaves on opens}

For an open \(W\) in a space \(X\), let \(P(W)\) be the free
abelian sheaf represented by \(W\). Concretely, sheafify the free
abelian presheaf on the representable presheaf of \(W\). What we
use is its natural representing property:

\[
\Hom(P(W),\mathcal M)\cong\Gamma(W,\mathcal M).
\]

An inclusion \(V\subseteq W\) induces \(P(V)\to P(W)\),
and \(P(X)\cong\underline{\Z}_X\). For a commutative ring \(R\),
write \(P_R(W)\) for the corresponding free sheaf of \(R\)-modules.
It represents sections in that category, with
\(P_R(X)\cong\underline R_X\).

\begin{lemma}\label{lem:projective-opens}
Suppose \(X\) is Hausdorff and totally disconnected. If \(W\) is
a disjoint union of compact-open subsets of \(X\), then \(P(W)\)
and \(P_R(W)\) are projective.\leanmark{projectivity}
\end{lemma}

\begin{proof}
To prove projectivity, it suffices to lift a section over \(W\)
through an epimorphism of sheaves. Epimorphisms are locally surjective,
so local lifts exist. On a compact open \(C\), choose finitely many
clopen sets refining a local lifting cover. This is possible because
\(C\) is compact Hausdorff and totally disconnected, hence
zero-dimensional. Subtract earlier clopen sets from later ones to
make the finite cover disjoint. The local lifts then glue, with no
compatibility conditions on overlaps, to a section on \(C\).
Do this on each compact-open component of \(W\) and glue again.
The representing property gives projectivity in both coefficient
categories.
\end{proof}

\subsection{The ordered complex and its exactness}

For the cover \(W_0,\ldots,W_r\), set

\[
P_k=\bigoplus_{0\leq i_0<\cdots<i_k\leq r}
P(W_{i_0}\cap\cdots\cap W_{i_k}).
\]

Deleting one index enlarges the intersection and induces a map of free
sheaves. Let the differential be the alternating sum of these maps.
The augmentation \(P_0\to\underline\Z_{U_r}\) is induced by the
inclusions of the charts. The usual face-deletion identities imply
that the differential squares to zero.

\begin{proposition}\label{prop:resolution}
The augmented complex

\[
0\longrightarrow P_r\longrightarrow\cdots\longrightarrow P_1
\longrightarrow P_0\longrightarrow\underline\Z_{U_r}
\longrightarrow0
\]

is a projective resolution. Replacing each free integral sheaf by
its \(R\)-module counterpart gives a projective resolution
\(P_\bullet^R\to\underline R_{U_r}\).\leanmark{resolution}
\end{proposition}

\begin{proof}
All terms are projective by \cref{lem:cells,lem:projective-opens}.
For exactness, take a stalk at \(x\in U_r\). The stalk of \(P(W)\)
is zero when \(x\notin W\); if \(x\in V\subseteq W\), the map
\(P(V)_x\to P(W)_x\) is an isomorphism. These assertions follow
from the representing property, or from the stalk--skyscraper
adjunction.

Let \(a\) be the least index with \(x\in W_a\). The surviving
faces form the simplex of charts containing \(x\), and its
augmented chain complex contracts to the vertex \(a\). Explicitly,
prepend \(a\) to an increasing face whose first vertex is larger
than \(a\), using the inverse of the stalk inclusion; on a face
starting at \(a\), set the contraction to zero. Faces starting below
\(a\) have zero stalk. In the alternating boundary formula, deleting
the inserted vertex gives the original face and all other terms
cancel in pairs. Thus \(dh+hd=1\), with the augmentation section
included in degree zero. The augmented stalk complex is exact.
Exactness of sheaves is detected on stalks. The same contraction works
for the module sheaves.
\end{proof}

Applying \(\Hom(-,\mathcal M)\) now computes ordinary sheaf
cohomology, since

\[
H^q(U_r;\mathcal M)
=\Ext^q_{\Sh(U_r;\Ab)}(\underline\Z,\mathcal M).
\]

This also identifies the cochain complex explicitly: its degree-\(k\)
term is the product of the section groups on intersections of \(k+1\)
charts, and its differential is the alternating sum of restrictions.
In particular, with constant \(\F\) coefficients, these sections are
locally constant \(\F\)-valued functions. Because there are no faces
of degree greater than \(r\), we have proved

\[
H^q(U_r;\mathcal M)=0\qquad(q>r)
\]

for every abelian coefficient sheaf.\leanmark{cohomology}

The remaining question is whether a particular cochain in the last
possible degree is a boundary. We first identify that cochain as a
power of a degree-one class.

\section{The class and the cup-product formula}
\label{sec:cup}

\subsection{Transition functions}

On \(W_i\cap W_j\), define

\[
z_{ij}=s_i+s_j\in\F,
\]

using any representative of the orbit. This is well defined: flipping
both bits adds \(1+1=0\). It is locally constant, as is visible on
the compact-open cells. For three indices,

\[
z_{ij}+z_{jk}=z_{ik}.
\]

Thus \(z=(z_{ij})_{i<j}\) is a degree-one cocycle in the complex
of \cref{sec:resolution}. Let

\[
\eta_r=[z]\in H^1(U_r;\F).
\]

Geometrically, \(z_{ij}\) records whether the gauge choices on the
two charts agree or differ by the involution. This observation
motivates the cocycle; the proof uses the displayed functions and
their identities directly.\leanmark{cocycle}

\subsection{Why multiplication is computed by adjacent differences}

\begin{proposition}\label{prop:cup-formula}
For \(1\leq k\leq r\), the class \(\eta_r^k\) is represented on
the face \([i_0,\ldots,i_k]\) by

\[
z_{i_0i_1}z_{i_1i_2}\cdots z_{i_{k-1}i_k}.
\]

In particular, on the full intersection and pulled back to configurations
in which every coordinate is finite, \(\eta_r^r\) is represented by

\begin{equation}\label{eq:top-cochain}
F_r((a_i,s_i)_{i=0}^r)=\prod_{j=0}^{r-1}(s_j+s_{j+1}).
\end{equation}
\end{proposition}

The familiar ordered \v{C}ech cup formula suggests this answer. We give
a chain-level proof, which also fixes its relation to ordinary sheaf
cohomology and to the operation used in the formalization.

\begin{proof}
Use the resolution \(P_\bullet^{\F}\) in sheaves of \(\F\)-modules.
Its graded Ext group

\[
\Ext^*_{\Sh(U_r;\F\text{-}\mathrm{Mod})}
(\underline\F,\underline\F)
\]

has the Yoneda product: multiplication is composition of the
corresponding morphisms in the derived category. This is the standard
realization of the cup algebra of derived global sections.

We compare this description with the ordinary cohomology already
computed. The integral and module resolutions are each independently
projective and exact. Their Hom complexes, with constant \(\F\)
as target, identify with the same complex of sections on chart
intersections. The identifications commute with restrictions and
hence with differentials. At degree zero the rank-one generator maps
to the constant function \(1\). This section-induced comparison
identifies the module Ext description with ordinary
\(H^*(U_r;\F)\), with its usual cup multiplication.\leanmark{comparison}

For a direct product calculation, write
\(d_k:P_{k+1}^{\F}\to P_k^{\F}\), and define

\[
L_k:P_{k+1}^{\F}\longrightarrow P_k^{\F}
\]

on a face \([i_0,\ldots,i_{k+1}]\) as follows: multiply by
\(z_{i_0i_1}\), delete the first vertex, and map into the summand
indexed by \([i_1,\ldots,i_{k+1}]\). A locally constant scalar
acts on sheaf sections by gluing the constant scalar actions on its
open fibers, so this construction defines a sheaf morphism.

The cocycle identity gives the chain-lift relation

\[
L_k\circ d_{k+1}=-d_k\circ L_{k+1}.
\]

Indeed, terms deleting a vertex after the first two cancel by the
face identities. The remaining terms combine by
\(z_{i_1i_2}+z_{i_0i_2}=z_{i_0i_1}\). In characteristic two the
minus sign is the same as a plus sign. If
\(\epsilon:P_0^{\F}\to\underline\F\) is the augmentation,
then \(\epsilon\circ L_0\), under the section comparison, is
precisely \(z\). Thus \(L\) lifts the degree-one cocycle to a
shifted map of the resolution.

Set \(c_0=\epsilon\) and \(c_{k+1}=c_k\circ L_k\).
The chain identity shows that these are cocycles. To see which classes
they represent, recall that the augmentation is inverted when
passing from the resolution to its resolved object in the derived
category. In composing two morphisms represented through the
resolution, the middle augmentation and its inverse cancel.
Consequently composing with the lift \(L\) computes multiplication
by \(\eta_r\). Induction shows that \(c_k\) represents
\(\eta_r^k\) for \(k\geq1\).

Finally, each application of \(L\) removes the current first vertex
and multiplies by its transition to the next one. On a face of length
\(k+1\), the composite therefore multiplies successively by
\(z_{i_0i_1},z_{i_1i_2},\ldots,z_{i_{k-1}i_k}\). This is the
claimed formula.\leanmark{cup}
\end{proof}

The product formula gives a useful simplification: although the space
has points at infinity, the top intersection has none. Its pulled-back
cochain \(F_r\) is just a function on a product of discrete sets.
The points at infinity return only when we ask whether this function
is a sum of restrictions of continuous functions on the intersections
one degree lower.
That is where the nonvanishing argument begins.

\section{Why the top cup power does not vanish}\label{sec:nonvanishing}

The cup-product calculation has reduced nonvanishing to a question about
functions.  We must show that the product in \eqref{eq:top-cochain} cannot be
the restriction of a coboundary.  Two features of the construction work
together: approaching infinity forces eventual constancy in a coordinate,
and the strictly increasing coordinate cardinalities let us use this
constancy simultaneously for every choice of the other coordinates.

Recall that \(D_i=A_i\times\F\), with \(|A_i|=\aleph_i\).  Write

\[
C_r=\prod_{i=0}^{r}D_i,
\qquad
\sigma_r((a_i,s_i)_{i=0}^{r})=(a_i,s_i+1)_{i=0}^{r}.
\]

Thus \(C_r\) consists of the all-finite configurations, before passing to
the quotient.  The pullback of the top product cochain is

\[
F_r((a_i,s_i)_{i=0}^{r})
=\prod_{j=0}^{r-1}(s_j+s_{j+1}).
\]

We use the empty-product convention \(F_0=1\).  Throughout this section,
all function values and sums are in \(\F\).

The function \(F_r\) depends only on the bits, not on the base points
\(a_i\).  For instance, it takes the value one whenever consecutive bits
are different.  This observation alone says nothing about cohomological
nonvanishing: a nonzero function can still be a coboundary.  The role of the
large base sets is to restrict which invariant functions can occur in a
boundary, through the continuity requirements at the points at infinity.
We exploit those requirements without trying to describe all possible
continuous functions on the intersections.

\subsection{What a boundary would imply}

Call a function \(f:C_r\to\F\) \emph{almost constant in coordinate \(i\)}
if, after fixing every other coordinate, its restriction to \(D_i\) is
constant outside a finite set.  Neither the exceptional set nor the
eventual constant is required to be independent of the fixed coordinates.

\begin{lemma}\label{lem:boundary-necessary}
Suppose \(r\geq1\).  If the top product cochain is a coboundary, then

\[
F_r=\sum_{i=0}^{r}f_i,
\]

where every \(f_i:C_r\to\F\) is invariant under \(\sigma_r\) and almost
constant in coordinate \(i\).
\end{lemma}

\begin{proof}
An ordered cochain one degree below the top has a component on each
intersection

\[
W_{\widehat i}=\bigcap_{j\neq i}W_j.
\]

These components are locally constant \(\F\)-valued functions.
The top differential is their sum after restriction to the full
intersection: its alternating signs disappear in characteristic two.
Pulling the components back to \(C_r\) gives the functions \(f_i\).
Because they are pulled back from the orbit space, they are invariant
under the simultaneous flip.

Fix every coordinate except \(i\).  Allowing coordinate \(i\) to vary
through its entire one-point compactification defines a continuous map

\[
D_i^+\longrightarrow W_{\widehat i}.
\]

This map is defined even at infinity.  Indeed, \(r\geq1\), so at least one
other coordinate remains finite; the removed all-infinite orbit is never
encountered.  Compose with the component function on \(W_{\widehat i}\).
Since \(\F\) is discrete, continuity at infinity says that this composite
equals its value at infinity on a cofinite subset of \(D_i\).  This is
exactly the required almost constancy.  Notice that this concerns the
whole doubled set \(D_i=A_i\times\F\).  Outside a finite set, both bit
values therefore have the same eventual value.  Two unrelated eventual
values on the separate bit-zero and bit-one subsets would not suffice
for the argument below.
\end{proof}

Only this necessary condition is used.  In particular, we do not claim
that an arbitrary family of coordinatewise almost-constant functions
extends to jointly continuous functions on the deleted intersections.
\leanmark{boundary}

\subsection{One choice avoiding all exceptional sets}

The exceptional set just obtained can depend on many other coordinates.
The following elementary counting argument supplies a single choice
that avoids all of them.

\begin{lemma}[Uniform thinning]\label{lem:uniform-thinning}
Let \(P\) be infinite and \(|P|<|A|\).  For each \(y\in P\), let
\(N_y\subseteq A\times\F\) be finite.  There is an \(a\in A\) such that
both \((a,0)\) and \((a,1)\) lie outside every \(N_y\).
\end{lemma}

\begin{proof}
Let \(\pi:A\times\F\to A\) be the first-coordinate projection.  Since
each \(\pi(N_y)\) is finite,

\[
\left|\bigcup_{y\in P}\pi(N_y)\right|
\leq |P|\cdot\aleph_0
=|P|
<|A|.
\]

Choose \(a\) outside this union.  No point of either form \((a,0)\) or
\((a,1)\) can belong to any of the exceptional sets.
\end{proof}

For the application, the index set is \(C_{r-1}\), whose cardinality is
\(\aleph_{r-1}\), while the last base set has cardinality \(\aleph_r\).
This strict inequality is the reason for increasing the coordinate sizes.
The argument uses no uniform bound on the sizes of the finite exceptional
sets, and no common eventual value.  It is simply a bound on a union of
finite sets, so no regularity assumption on the larger cardinal is being
used.

Choosing a point separately for each configuration would be much weaker.
It would give a base point depending on the remaining coordinates, and
then taking the associated slices could destroy almost constancy in
those coordinates.  The uniform choice instead gives two fixed slices
of every function.  Their algebraic and continuity properties can then
be compared directly.  This is the only place where the strict inequality
between successive coordinate sizes is used; the remaining reduction is
elementary linear algebra.  \leanmark{thinning}

\subsection{Removing one coordinate by a finite difference}

\begin{proposition}\label{prop:nondecomposition}
For every \(r\geq1\), there is no decomposition

\[
F_r=\sum_{i=0}^{r}f_i
\]

in which \(f_i\) is invariant under \(\sigma_r\) and almost constant in
coordinate \(i\).
\end{proposition}

\begin{proof}
We establish a reduction in degree and then prove the two-coordinate
base case separately.  Suppose that such a decomposition has been given.
For each \(y\in C_{r-1}\), choose an eventual value \(c_y\) and a finite
exceptional set \(N_y\subseteq D_r\) for the last-coordinate function
\(z\mapsto f_r(y,z)\).  Lemma~\ref{lem:uniform-thinning} supplies
\(a\in A_r\) avoiding all these exceptions.

For any function \(f:C_r\to\F\), define its two-bit difference by

\[
(\Delta_a f)(y)=f(y,(a,0))+f(y,(a,1)),
\qquad y\in C_{r-1}.
\]

For our chosen \(a\), both values of the last summand equal \(c_y\).
Consequently, \(\Delta_a f_r=0\).  In contrast, the product function
loses exactly its last factor:

\[
\begin{aligned}
(\Delta_a F_r)(y)
&=F_{r-1}(y)\bigl(s_{r-1}+(s_{r-1}+1)\bigr)\\
&=F_{r-1}(y).
\end{aligned}
\]

For \(i<r\), the function \(\Delta_a f_i\) remains almost constant in
coordinate \(i\).  Fix the other coordinates and take the union of the
finite exceptional sets for its two slices.  Outside this union, the sum
of their two eventual values is constant.

The difference also preserves invariance.  If \(f\) is invariant under
the simultaneous flip, then

\[
\begin{aligned}
(\Delta_a f)(\sigma_{r-1}y)
&=f(\sigma_{r-1}y,(a,0))+f(\sigma_{r-1}y,(a,1))\\
&=f(y,(a,1))+f(y,(a,0))\\
&=(\Delta_a f)(y).
\end{aligned}
\]

Here it matters that the two points have the same base coordinate:
the flip exchanges the two summands.  Thus, for \(r\geq2\), applying
\(\Delta_a\) gives

\[
F_{r-1}=\sum_{i=0}^{r-1}\Delta_a f_i,
\]

with all the required properties in one lower positive degree.  This
contradicts the induction hypothesis once the case \(r=1\) is known.

In this reduction, the largest coordinate and the summand controlled by
that coordinate disappear together.  Each surviving summand retains its
own almost-constancy condition, while the target function becomes the
same adjacent-bit product on the shorter configuration space.  The
operation sums only two values, so no infinite sums or convergence
questions enter the reduction.

For that base case, the reduction gives \(\Delta_a f_0=1\).  Retain the
single slice

\[
g(x)=f_0(x,(a,0)),
\qquad x\in D_0.
\]

It is almost constant because \(f_0\) has that property in coordinate
zero.  Invariance of \(f_0\), applied to \((x,(a,1))\), gives
\(g(\sigma_0x)=f_0(x,(a,1))\).  Hence

\[
g(x)+g(\sigma_0x)=1
\qquad\text{for every }x\in D_0.
\]

Choose a finite exceptional set \(N\subseteq D_0\) for \(g\), with
eventual value \(c\).  The set \(A_0\) is infinite, so it contains
\(a_0\) outside the first-coordinate projection of \(N\).  Both
\((a_0,0)\) and \((a_0,1)\) then lie outside \(N\).  Substituting the
first of these points into the displayed identity gives \(c+c=1\),
contradicting \(c+c=0\).  This proves the base case and the induction.\leanmark{difference}
\end{proof}

\begin{remark}
Two tempting shortcuts would lose the argument.  First, the single slice
\(g\) need not be invariant; the proof uses the relation between \(g\)
and its flipped value instead.  Second, one cannot continue a bare
nondecomposition induction down to degree zero.  The function \(F_0=1\)
is itself invariant and almost constant, so the corresponding
nondecomposition assertion is false.  The base case succeeds because it
retains the extra information that this constant function is the flip
difference of an almost-constant slice.
\end{remark}

\subsection{Returning to cohomology}

The top cochain in \eqref{eq:top-cochain} represents \(\eta_r^r\), by the
chain-lift calculation.  If this class were zero, the projective
resolution would express that cochain as a coboundary.
Lemma~\ref{lem:boundary-necessary} would then produce the decomposition
forbidden by Proposition~\ref{prop:nondecomposition}.  We conclude that

\[
\eta_r^r\neq0.
\]

The same resolution has length \(r\), so it also gives
\(H^{r+1}(U_r;\F)=0\), and therefore \(\eta_r^{r+1}=0\).  Combined with
the cardinality and weight of \(U_r\), these statements prove the
existence assertions of Theorem~\ref{thm:main}.  \leanmark{nonzero}

\section{The universal lower bound}\label{sec:optimality}

The examples constructed above attain the least possible cardinality and
weight. We prove this by an estimate for arbitrary abelian coefficient
sheaves, independent of the quotient construction. The relevant invariant
is the number of compact opens needed to cover a space.

Recall that \(P(W)\) denotes the integral free sheaf represented by an open
subset \(W\subseteq X\), so that

\[
\Hom(P(W),\mathcal M)\cong\Gamma(W,\mathcal M).
\]

In particular, \(P(X)\cong\underline{\Z}_X\). Throughout this section,
projective dimension is understood in its Ext-vanishing sense:
\(\pd(P)\leq n\) means that \(\Ext^d(P,M)=0\) for every coefficient object
\(M\) and every \(d>n\). The category of abelian sheaves has enough
injectives, so dimension shifting in the coefficient variable is available.
The indexing begins with countable covers: they already give projective
free-open sheaves, rather than merely a dimension bound of one. Each passage
to the next infinite cardinal costs at most one further degree. Keeping
this starting point explicit will give precisely the lower bound
\(\aleph_r\) in degree \(r\).

\subsection{Continuous transfinite extensions}

The homological input is a closure property whose proof is worth recording.
Continuity of the filtration is essential to its use below.

\begin{lemma}\label{lem:transfinite-dimension}
Let \(\mathcal A\) be an abelian category with enough injectives. Suppose
\((P_\alpha)_{\alpha\leq\lambda}\) is a well-ordered filtration with
\(P_0=0\), monomorphic successor maps, and

\[
P_\delta\cong\mathop{\mathrm{colim}}_{\alpha<\delta}P_\alpha
\]

at every nonzero limit ordinal \(\delta\leq\lambda\). If each successive
quotient \(Q_\alpha=P_{\alpha+1}/P_\alpha\) has projective dimension at
most \(n\), then \(\pd(P_\lambda)\leq n\).
\end{lemma}

\begin{proof}
First fix \(Y\in\mathcal A\) and assume only that
\(\Ext^1(Q_\alpha,Y)=0\) for all \(\alpha<\lambda\). Choose an injective
presentation

\[
0\longrightarrow Y\longrightarrow I\xrightarrow{p}C\longrightarrow0.
\]

The long exact Ext sequence shows that \(\Ext^1(Q,Y)=0\) precisely when
every map \(Q\to C\) lifts through \(p\). We will prove this lifting
property for \(P_\lambda\).

Given \(f:P_\lambda\to C\), denote its restriction to \(P_\alpha\) by
\(f_\alpha\). Construct compatible maps \(g_\alpha:P_\alpha\to I\)
satisfying \(p g_\alpha=f_\alpha\). There is a unique choice at stage zero.
At a successor stage, injectivity of \(I\) extends \(g_\alpha\) to a map
\(h:P_{\alpha+1}\to I\). This extension may not lift \(f_{\alpha+1}\),
but the discrepancy \(f_{\alpha+1}-ph\) vanishes on \(P_\alpha\).
Writing \(\pi_\alpha:P_{\alpha+1}\to Q_\alpha\) for the quotient map,
there is therefore a map \(b:Q_\alpha\to C\) such that

\[
f_{\alpha+1}-ph=b\pi_\alpha.
\]

By the assumed Ext vanishing, choose \(\widetilde b:Q_\alpha\to I\)
with \(p\widetilde b=b\). Then
\(g_{\alpha+1}=h+\widetilde b\pi_\alpha\) is both an extension of
\(g_\alpha\) and a lift of \(f_{\alpha+1}\). At a limit stage, the
compatible earlier lifts induce a map from the colimit. Their identities
with the restrictions of \(f\) persist by the colimit's universal property.
This constructs \(g_\lambda\), proving \(\Ext^1(P_\lambda,Y)=0\).

Now fix any \(d>n\) and any coefficient object \(Y\). Choose successive
injective presentations

\[
0\longrightarrow Y_j\longrightarrow I_j\longrightarrow Y_{j+1}
\longrightarrow0,
\qquad Y_0=Y,
\qquad 0\leq j<d-1.
\]

Dimension shifting gives, for every object \(Q\),

\[
\Ext^d(Q,Y)\cong\Ext^1(Q,Y_{d-1}).
\]

The same coefficient object \(Y_{d-1}\) works for all successive
quotients. Their dimension bounds make the right-hand side vanish, so the
degree-one argument applies. Shifting back gives
\(\Ext^d(P_\lambda,Y)=0\). Since \(d>n\) and \(Y\) were arbitrary,
the required projective-dimension bound follows.\leanmark{transfinite}
\end{proof}

\begin{remark}
The proof constructs compatible lifts; it does not commute Ext with an
arbitrary filtered colimit. In particular, the closure statement requires
no separate assumption that arbitrary filtered colimits are exact.
\end{remark}

\subsection{Counting compact opens}

\begin{proposition}\label{prop:compact-open-dimension}
Let \(X\) be Hausdorff and totally disconnected. If an open subset
\(W\subseteq X\) is a union of at most \(\aleph_n\) compact-open subsets
of \(X\), then \(\pd(P(W))\leq n\).
\end{proposition}

\begin{proof}
We induct on \(n\). For \(n=0\), enumerate the given compact opens as
\(K_0,K_1,\ldots\), allowing empty terms, and set

\[
L_j=K_j\setminus\bigcup_{i<j}K_i.
\]

The earlier union is finite and compact, hence closed in the Hausdorff
space \(X\); consequently \(L_j\) is open. It is also closed in the
compact set \(K_j\), since the earlier union is open. Thus the \(L_j\)
are compact open, pairwise disjoint, and have union \(W\).
Lemma~\ref{lem:projective-opens} makes \(P(W)\) projective. This includes
finite and empty covers.
Concretely, sections through a sheaf epimorphism lift separately on every
\(L_j\), and the lifts glue without compatibility conditions because the
pieces are disjoint. This explains why no finiteness condition on the
disjoint family is needed. Local compactness of the ambient space is not
required for this step or for the proposition.

Assume the assertion for \(n\). Pad a family of at most
\(\aleph_{n+1}\) compact opens with empty sets to index it by the initial
ordinal \(\kappa=\omega_{n+1}\), and put

\[
W_\alpha=\bigcup_{\beta<\alpha}K_\beta
\quad(\alpha\leq\kappa),
\qquad W_\kappa=W.
\]

These opens give a continuous filtration of free sheaves. Indeed,
inclusions induce monomorphisms: on stalks the free-open sheaf is \(\Z\)
inside its representing open and zero outside, and the nonzero inclusion
maps are identities. Moreover, the free sheaf of a directed union is the
colimit of the free sheaves of its members. By the representing property,
this assertion says exactly that compatible sections on the members glue
uniquely to a section on their union. Hence \(P(W_0)=0\), continuity holds
at limit ordinals, and the final object is \(P(W)\).

For two opens \(V,T\), the free-open Mayer--Vietoris sequence is

\[
0\longrightarrow P(V\cap T)
\longrightarrow P(V)\oplus P(T)
\longrightarrow P(V\cup T)\longrightarrow0.
\]

The first map is the signed pair of inclusions, and the second adds the
inclusions into the union. Exactness can be checked on stalks; on the
intersection the maps are \(u\mapsto(u,-u)\) and \((v,t)\mapsto v+t\).
It follows, by taking the quotient by \(P(V)\), that the successive
quotient \(Q_\alpha=P(W_{\alpha+1})/P(W_\alpha)\) fits into

\[
0\longrightarrow P(W_\alpha\cap K_\alpha)
\longrightarrow P(K_\alpha)
\longrightarrow Q_\alpha\longrightarrow0.
\]

For every \(\alpha<\kappa\), the initial-ordinal property gives
\(|\alpha|<\aleph_{n+1}\), hence \(|\alpha|\leq\aleph_n\). Since

\[
W_\alpha\cap K_\alpha
=\bigcup_{\beta<\alpha}(K_\beta\cap K_\alpha),
\]

the induction hypothesis bounds the left-hand term's projective dimension
by \(n\). The middle term is projective by
Lemma~\ref{lem:projective-opens}. For \(d>n+1\), the long exact Ext
sequence identifies \(\Ext^d(Q_\alpha,M)\) with
\(\Ext^{d-1}(P(W_\alpha\cap K_\alpha),M)\), which vanishes.
Thus \(\pd(Q_\alpha)\leq n+1\).
Lemma~\ref{lem:transfinite-dimension} completes the induction.
Only the size of each proper initial segment of \(\omega_{n+1}\) enters
the argument. There is no appeal to the continuum hypothesis or to any
extra cardinal-arithmetic assumption.\leanmark{cover-bound}
\end{proof}

Taking \(W=X\), the identification \(P(X)\cong\underline{\Z}_X\)
immediately gives the following consequence.

\begin{corollary}\label{cor:cover-vanishing}
If a Hausdorff totally disconnected space \(X\) has a compact-open cover
of cardinality at most \(\aleph_n\), then
\(H^d(X;\mathcal M)=0\) for every abelian sheaf \(\mathcal M\) and
every \(d>n\).
\end{corollary}

\subsection{Cardinality and weight}

\begin{corollary}\label{cor:universal-minimality}
Let \(V\) be locally compact, Hausdorff, and totally disconnected, let
\(r\geq1\), and let \(\mathcal M\) be any abelian sheaf on \(V\).
Then

\[
H^r(V;\mathcal M)\neq0
\quad\Longrightarrow\quad
|V|\geq\aleph_r
\quad\text{and}\quad
w(V)\geq\aleph_r.
\]

\end{corollary}

\begin{proof}
Such a space has a basis of compact opens. To recall why, given a point
and an open neighborhood, choose a compact neighborhood inside it. This
compact Hausdorff totally disconnected subspace is zero-dimensional, so
it contains a clopen neighborhood of the point lying inside its own
interior in \(V\). That smaller neighborhood is both open in \(V\) and
compact. Thus compact opens refine every neighborhood.

Choosing one compact-open
neighborhood of each point gives a cover with at most \(|V|\) members.
Thus \(|V|\leq\aleph_n\) implies vanishing above degree \(n\) by
Corollary~\ref{cor:cover-vanishing}.

For weight, choose a basis \(\mathcal B\) of cardinality \(w(V)\).
Every compact open is a finite union of basis elements: cover it by basis
neighborhoods contained in it and use compactness. Choosing one such
finite description for each compact open gives

\[
|\operatorname{CompactOpens}(V)|
\leq\max(\aleph_0,w(V)).
\]

Indeed, identical finite descriptions have identical unions, so this
choice is injective into the set of finite subsets of \(\mathcal B\).
For an infinite basis, that set has the same cardinality as the basis.
All compact opens together cover \(V\). Therefore \(w(V)\leq\aleph_n\)
also forces vanishing above degree \(n\); the displayed maximum handles
finite or empty bases as well.

Finally, a cardinal strictly below \(\aleph_r\), for positive finite
\(r\), is at most \(\aleph_{r-1}\). Either \(|V|<\aleph_r\) or
\(w(V)<\aleph_r\) would therefore imply \(H^r(V;\mathcal M)=0\).
Taking contrapositives proves both inequalities.\leanmark{minimality}
\end{proof}

A nonzero product of \(r\) degree-one classes is a nonzero degree-\(r\)
cohomology class, so these inequalities apply to every space considered in
Aoki's question. The examples of Theorem~\ref{thm:main} attain both bounds
simultaneously. This proves the asserted optimality for cardinality and
weight.

\section{What the formalization checks}
\label{sec:formalization}

The proof, formalization, and exposition were developed with AI assistance.
The accompanying development formalizes the spaces, their topological
properties, both exact size calculations, the finite projective
resolutions, the cocycle, the cup-product calculation, the nonvanishing
argument, and the universal lower bounds. The assembled theorem is about
actual spaces and ordinary sheaf cohomology, rather than an auxiliary
quotient defined only to encode the combinatorial obstruction.

There are three points of translation between the exposition and the
code. First, the integral free sheaf on the whole space is identified
with the constant integral object in Mathlib's definition of ordinary
sheaf cohomology. Second, the cup operation is implemented through the
constant-module Yoneda algebra and the proved comparison of section
complexes from \cref{sec:cup}. The code proves the chain-lift formula;
it does not compare this construction with an independently defined
Mathlib sheaf cup-product operation. Third, although the gauge choices
suggest interpreting \(\eta_r\) as the class of a double cover, a
torsor-classification equivalence is not needed and is not formalized.

The code uses Boolean bits and the field with two elements. Its recursive
configuration type indexed by \(r\) has \(r+1\) coordinates, and its
alternating function has \(r\) factors. The final theorem uses the same
positive degree \(r\) as \cref{thm:main}; some inductive helpers use
an index one smaller.

\subsection{Reproducing the verification}

The project pins Lean v4.31.0 and Mathlib v4.31.0, with Mathlib commit
\begin{center}
\small\nolinkurl{fabf563a7c95a166b8d7b6efca11c8b4dc9d911f}.
\end{center}
The transitive dependency commits are recorded in the manifest. From the
project's \texttt{lean} directory, run
\begin{verbatim}
lake exe cache get
lake build
lake env lean AxiomAudit.lean
\end{verbatim}
The aggregate \texttt{Aoki.lean} imports 64 local modules. The recorded
build succeeds, and the axiom audit for the principal statements and
bridges reports only propositional extensionality, classical choice,
and quotient soundness. The sources contain no proof placeholders or
additional mathematical axioms. The output is retained in
\href{../../../lean/verification.txt}{\texttt{lean/verification.txt}}.

\Cref{app:lean} is a source concordance, not an additional mathematical
assumption: each entry identifies definitions and proved declarations
used in the proof above. File links are relative to the distributed
project, and printed line numbers locate the declarations even in a
viewer that does not open Lean source files.

\appendix
\section{Lean source concordance}
\label{app:lean}

The following entries give the precise source locations for the markers
in the text. All paths are relative to the \texttt{lean} directory;
line numbers refer to the accompanying source snapshot. Each file name
is a link to that file in the distributed project. The declarations are
grouped by their mathematical role so that a reader can follow one part
of the argument without reading the entire development.

\leanentry{main}{The assembled theorem}
\begin{leansources}
\item \href{../../../lean/Aoki/Main.lean}{\nolinkurl{Aoki/Main.lean}}, line 31:\\*
\nolinkurl{aoki_question_2_14}.
\item \href{../../../lean/Aoki/Main.lean}{\nolinkurl{Aoki/Main.lean}}, line 47:\\*
\nolinkurl{aoki_question_2_14_minimality}.
\item \href{../../../lean/Aoki/Main.lean}{\nolinkurl{Aoki/Main.lean}}, line 56:\\*
\nolinkurl{aoki_example_cohomology_above}.
\end{leansources}

\leanentry{construction}{The punctured involution quotient}
\begin{leansources}
\item \href{../../../lean/Aoki/Topology/Quotient.lean}{\nolinkurl{Aoki/Topology/Quotient.lean}}, line 19:\\*
\nolinkurl{CompactProduct}.
\item \href{../../../lean/Aoki/Topology/Quotient.lean}{\nolinkurl{Aoki/Topology/Quotient.lean}}, line 138:\\*
\nolinkurl{CompactQuotient}.
\item \href{../../../lean/Aoki/Topology/Quotient.lean}{\nolinkurl{Aoki/Topology/Quotient.lean}}, line 163:\\*
\nolinkurl{U}.
\item \href{../../../lean/Aoki/Topology/Quotient.lean}{\nolinkurl{Aoki/Topology/Quotient.lean}}, line 73:\\*
\nolinkurl{flip_eq_self_iff}.
\item \href{../../../lean/Aoki/Topology/Quotient.lean}{\nolinkurl{Aoki/Topology/Quotient.lean}}, line 196:\\*
\nolinkurl{isOpenMap_quotientMap}.
\end{leansources}

\leanentry{charts}{Gauge charts and compact-open cells}
\begin{leansources}
\item \href{../../../lean/Aoki/Sheaf/ChartProjectivity.lean}{\nolinkurl{Aoki/Sheaf/ChartProjectivity.lean}}, line 37:\\*
\nolinkurl{iSup_chartOpen}.
\item \href{../../../lean/Aoki/Topology/Charts.lean}{\nolinkurl{Aoki/Topology/Charts.lean}}, line 86:\\*
\nolinkurl{sliceHomeomorph}.
\item \href{../../../lean/Aoki/Topology/Charts.lean}{\nolinkurl{Aoki/Topology/Charts.lean}}, line 126:\\*
\nolinkurl{chartProductHomeomorph}.
\item \href{../../../lean/Aoki/Topology/Intersections.lean}{\nolinkurl{Aoki/Topology/Intersections.lean}}, line 14:\\*
\nolinkurl{GaugeLabel}.
\item \href{../../../lean/Aoki/Topology/Intersections.lean}{\nolinkurl{Aoki/Topology/Intersections.lean}}, line 151:\\*
\nolinkurl{cellOpen}.
\item \href{../../../lean/Aoki/Topology/Intersections.lean}{\nolinkurl{Aoki/Topology/Intersections.lean}}, line 84:\\*
\nolinkurl{cell_pairwiseDisjoint}.
\item \href{../../../lean/Aoki/Topology/Intersections.lean}{\nolinkurl{Aoki/Topology/Intersections.lean}}, line 92:\\*
\nolinkurl{cells_cover_intersection}.
\item \href{../../../lean/Aoki/Topology/Intersections.lean}{\nolinkurl{Aoki/Topology/Intersections.lean}}, line 126:\\*
\nolinkurl{isClopen_cell}.
\item \href{../../../lean/Aoki/Topology/Intersections.lean}{\nolinkurl{Aoki/Topology/Intersections.lean}}, line 130:\\*
\nolinkurl{isCompact_cell}.
\end{leansources}

\leanentry{size}{Cardinality and weight}
\begin{leansources}
\item \href{../../../lean/Aoki/Topology/CompactWeight.lean}{\nolinkurl{Aoki/Topology/CompactWeight.lean}}, line 18:\\*
\nolinkurl{weight_le_max_mk}.
\item \href{../../../lean/Aoki/Topology/Axis.lean}{\nolinkurl{Aoki/Topology/Axis.lean}}, line 37:\\*
\nolinkurl{isEmbedding_axis}.
\item \href{../../../lean/Aoki/Topology/GeometrySize.lean}{\nolinkurl{Aoki/Topology/GeometrySize.lean}}, line 52:\\*
\nolinkurl{mk_U}.
\item \href{../../../lean/Aoki/Topology/GeometrySize.lean}{\nolinkurl{Aoki/Topology/GeometrySize.lean}}, line 74:\\*
\nolinkurl{weight_U}.
\item \href{../../../lean/Aoki/Topology/Weight.lean}{\nolinkurl{Aoki/Topology/Weight.lean}}, line 26:\\*
\nolinkurl{weight}.
\end{leansources}

\leanentry{projectivity}{Represented free sheaves and projectivity}
\begin{leansources}
\item \href{../../../lean/Aoki/Sheaf/FreeOpen.lean}{\nolinkurl{Aoki/Sheaf/FreeOpen.lean}}, line 33:\\*
\nolinkurl{freeOpenHomEquiv}.
\item \href{../../../lean/Aoki/Sheaf/FreeOpen.lean}{\nolinkurl{Aoki/Sheaf/FreeOpen.lean}}, line 97:\\*
\nolinkurl{freeOpenTopIso}.
\item \href{../../../lean/Aoki/Sheaf/ModuleFreeOpen.lean}{\nolinkurl{Aoki/Sheaf/ModuleFreeOpen.lean}}, line 31:\\*
\nolinkurl{freeOpenModuleHomEquiv}.
\item \href{../../../lean/Aoki/Sheaf/ModuleFreeOpen.lean}{\nolinkurl{Aoki/Sheaf/ModuleFreeOpen.lean}}, line 97:\\*
\nolinkurl{freeOpenModuleTopIso}.
\item \href{../../../lean/Aoki/Topology/CompactOpenCovers.lean}{\nolinkurl{Aoki/Topology/CompactOpenCovers.lean}}, line 79:\\*
\nolinkurl{exists_compactOpen_disjoint_refinement}.
\item \href{../../../lean/Aoki/Sheaf/EpiSections.lean}{\nolinkurl{Aoki/Sheaf/EpiSections.lean}}, line 70:\\*
\nolinkurl{surjective_sections_of_disjoint_refinements}.
\item \href{../../../lean/Aoki/Sheaf/ProjectiveOpen.lean}{\nolinkurl{Aoki/Sheaf/ProjectiveOpen.lean}}, line 41:\\*
\nolinkurl{projective_freeOpen_disjoint_iSup}.
\item \href{../../../lean/Aoki/Sheaf/ModuleFreeOpen.lean}{\nolinkurl{Aoki/Sheaf/ModuleFreeOpen.lean}}, line 184:\\*
\nolinkurl{projective_freeOpenModule_disjoint_iSup}.
\item \href{../../../lean/Aoki/Sheaf/ChartProjectivity.lean}{\nolinkurl{Aoki/Sheaf/ChartProjectivity.lean}}, line 64:\\*
\nolinkurl{projective_freeOpen_chartIntersection}.
\item \href{../../../lean/Aoki/Sheaf/ModuleChartProjectivity.lean}{\nolinkurl{Aoki/Sheaf/ModuleChartProjectivity.lean}}, line 25:\\*
\nolinkurl{projective_freeOpenModule_chartIntersection}.
\end{leansources}

\leanentry{resolution}{The finite projective resolutions}
\begin{leansources}
\item \href{../../../lean/Aoki/Cech/OrderedComplex.lean}{\nolinkurl{Aoki/Cech/OrderedComplex.lean}}, line 161:\\*
\nolinkurl{orderedComplex}.
\item \href{../../../lean/Aoki/Cech/OrderedComplex.lean}{\nolinkurl{Aoki/Cech/OrderedComplex.lean}}, line 178:\\*
\nolinkurl{orderedTerm_isZero}.
\item \href{../../../lean/Aoki/Cech/OrderedCone.lean}{\nolinkurl{Aoki/Cech/OrderedCone.lean}}, line 194:\\*
\nolinkurl{orderedContraction_identity}.
\item \href{../../../lean/Aoki/Cech/AugmentedCone.lean}{\nolinkurl{Aoki/Cech/AugmentedCone.lean}}, line 70:\\*
\nolinkurl{orderedContraction_augmentation_identity}.
\item \href{../../../lean/Aoki/Cech/SheafResolution.lean}{\nolinkurl{Aoki/Cech/SheafResolution.lean}}, line 49:\\*
\nolinkurl{freeOpenProjectiveResolution}.
\item \href{../../../lean/Aoki/Cech/ModuleSheafResolution.lean}{\nolinkurl{Aoki/Cech/ModuleSheafResolution.lean}}, line 81:\\*
\nolinkurl{freeOpenModuleProjectiveResolution}.
\end{leansources}

\leanentry{cohomology}{Ordinary cohomology and upper vanishing}
\begin{leansources}
\item \href{../../../lean/Aoki/Cech/NonzeroClass.lean}{\nolinkurl{Aoki/Cech/NonzeroClass.lean}}, line 32:\\*
\nolinkurl{quotientIntegralResolution}.
\item \href{../../../lean/Aoki/Cech/NonzeroClass.lean}{\nolinkurl{Aoki/Cech/NonzeroClass.lean}}, line 70:\\*
\nolinkurl{quotient_sheaf_cohomology_eq_zero_above}.
\item \href{../../../lean/Aoki/Sheaf/LocallyConstant.lean}{\nolinkurl{Aoki/Sheaf/LocallyConstant.lean}}, line 112:\\*
\nolinkurl{constantSheafSectionAddEquiv}.
\end{leansources}

\leanentry{cocycle}{The degree-one transition cocycle}
\begin{leansources}
\item \href{../../../lean/Aoki/Cech/Cocycle.lean}{\nolinkurl{Aoki/Cech/Cocycle.lean}}, line 70:\\*
\nolinkurl{transitionOnOpen}.
\item \href{../../../lean/Aoki/Cech/Cocycle.lean}{\nolinkurl{Aoki/Cech/Cocycle.lean}}, line 155:\\*
\nolinkurl{transitionCochain}.
\item \href{../../../lean/Aoki/Cech/Cocycle.lean}{\nolinkurl{Aoki/Cech/Cocycle.lean}}, line 167:\\*
\nolinkurl{transitionCochain_cocycle}.
\item \href{../../../lean/Aoki/Cech/NonzeroClass.lean}{\nolinkurl{Aoki/Cech/NonzeroClass.lean}}, line 51:\\*
\nolinkurl{degreeOneClass}.
\item \href{../../../lean/Aoki/Cech/Cocycle.lean}{\nolinkurl{Aoki/Cech/Cocycle.lean}}, line 119:\\*
\nolinkurl{constantF2Sheaf}.
\end{leansources}

\leanentry{comparison}{The module-to-abelian comparison}
\begin{leansources}
\item \href{../../../lean/Aoki/Cech/ModuleComparison.lean}{\nolinkurl{Aoki/Cech/ModuleComparison.lean}}, line 127:\\*
\nolinkurl{moduleCechHomEquiv_d}.
\item \href{../../../lean/Aoki/Sheaf/ModuleConstantComparison.lean}{\nolinkurl{Aoki/Sheaf/ModuleConstantComparison.lean}}, line 125:\\*
\nolinkurl{freeOpenModuleTopSectionAddEquiv_generator}.
\item \href{../../../lean/Aoki/Cech/CupCohomology.lean}{\nolinkurl{Aoki/Cech/CupCohomology.lean}}, line 86:\\*
\nolinkurl{cupCohomologyAddEquiv}.
\item \href{../../../lean/Aoki/Cech/CupCohomology.lean}{\nolinkurl{Aoki/Cech/CupCohomology.lean}}, line 104:\\*
\nolinkurl{cupCohomologyAddEquiv_extMk}.
\item \href{../../../lean/Aoki/Cech/CupCohomology.lean}{\nolinkurl{Aoki/Cech/CupCohomology.lean}}, line 118:\\*
\nolinkurl{cupProduct}.
\item \href{../../../lean/Aoki/Cech/CupCohomology.lean}{\nolinkurl{Aoki/Cech/CupCohomology.lean}}, line 127:\\*
\nolinkurl{cupPower}.
\end{leansources}

\leanentry{cup}{The chain lift and cup powers}
\begin{leansources}
\item \href{../../../lean/Aoki/Sheaf/LocallyConstantScalar.lean}{\nolinkurl{Aoki/Sheaf/LocallyConstantScalar.lean}}, line 72:\\*
\nolinkurl{locallyConstantSMul}.
\item \href{../../../lean/Aoki/Cech/CupProduct.lean}{\nolinkurl{Aoki/Cech/CupProduct.lean}}, line 45:\\*
\nolinkurl{cupLift}.
\item \href{../../../lean/Aoki/Cech/CupProduct.lean}{\nolinkurl{Aoki/Cech/CupProduct.lean}}, line 197:\\*
\nolinkurl{cupLift_chain_signed}.
\item \href{../../../lean/Aoki/Cech/CupProduct.lean}{\nolinkurl{Aoki/Cech/CupProduct.lean}}, line 231:\\*
\(\iota\)\nolinkurl{_cupIterate}.
\item \href{../../../lean/Aoki/Homological/ExtMkComposition.lean}{\nolinkurl{Aoki/Homological/ExtMkComposition.lean}}, line 116:\\*
\nolinkurl{extMk_comp_of_lift}.
\item \href{../../../lean/Aoki/Cech/CupExt.lean}{\nolinkurl{Aoki/Cech/CupExt.lean}}, line 62:\\*
\nolinkurl{yonedaPower_cupExtClass}.
\item \href{../../../lean/Aoki/Cech/CupComparison.lean}{\nolinkurl{Aoki/Cech/CupComparison.lean}}, line 51:\\*
\nolinkurl{cupIterate_comparison}.
\item \href{../../../lean/Aoki/Cech/CupNonvanishing.lean}{\nolinkurl{Aoki/Cech/CupNonvanishing.lean}}, line 49:\\*
\nolinkurl{cupPower_degreeOneClass_top}.
\end{leansources}

\leanentry{boundary}{Necessary properties of a boundary}
\begin{leansources}
\item \href{../../../lean/Aoki/OnePoint.lean}{\nolinkurl{Aoki/OnePoint.lean}}, line 64:\\*
\nolinkurl{almostConstant_of_continuous_onePoint}.
\item \href{../../../lean/Aoki/Cech/BoundaryObstruction.lean}{\nolinkurl{Aoki/Cech/BoundaryObstruction.lean}}, line 142:\\*
\nolinkurl{invariant_deleted_section}.
\item \href{../../../lean/Aoki/Cech/BoundaryObstruction.lean}{\nolinkurl{Aoki/Cech/BoundaryObstruction.lean}}, line 152:\\*
\nolinkurl{along_deleted_section}.
\item \href{../../../lean/Aoki/Cech/CochainValues.lean}{\nolinkurl{Aoki/Cech/CochainValues.lean}}, line 68:\\*
\nolinkurl{cochainValue_differential}.
\end{leansources}

\leanentry{thinning}{Uniform thinning and cardinal growth}
\begin{leansources}
\item \href{../../../lean/Aoki/Thinning.lean}{\nolinkurl{Aoki/Thinning.lean}}, line 18:\\*
\nolinkurl{mk_exceptional_first_coordinates_le}.
\item \href{../../../lean/Aoki/Thinning.lean}{\nolinkurl{Aoki/Thinning.lean}}, line 57:\\*
\nolinkurl{exists_fresh_pair}.
\item \href{../../../lean/Aoki/Cardinals.lean}{\nolinkurl{Aoki/Cardinals.lean}}, line 45:\\*
\nolinkurl{mk_config}.
\item \href{../../../lean/Aoki/Cardinals.lean}{\nolinkurl{Aoki/Cardinals.lean}}, line 62:\\*
\nolinkurl{config_lt_next}.
\end{leansources}

\leanentry{difference}{Finite differences and nondecomposition}
\begin{leansources}
\item \href{../../../lean/Aoki/Functions.lean}{\nolinkurl{Aoki/Functions.lean}}, line 67:\\*
\nolinkurl{delta_alternating}.
\item \href{../../../lean/Aoki/Functions.lean}{\nolinkurl{Aoki/Functions.lean}}, line 73:\\*
\nolinkurl{invariant_delta}.
\item \href{../../../lean/Aoki/AlmostConstant.lean}{\nolinkurl{Aoki/AlmostConstant.lean}}, line 78:\\*
\nolinkurl{not_flip_sum_one}.
\item \href{../../../lean/Aoki/Nonvanishing.lean}{\nolinkurl{Aoki/Nonvanishing.lean}}, line 73:\\*
\nolinkurl{alternating_one_not_boundary}.
\item \href{../../../lean/Aoki/Nonvanishing.lean}{\nolinkurl{Aoki/Nonvanishing.lean}}, line 51:\\*
\nolinkurl{boundary_delta}.
\item \href{../../../lean/Aoki/Nonvanishing.lean}{\nolinkurl{Aoki/Nonvanishing.lean}}, line 120:\\*
\nolinkurl{aleph_nondecomposition}.
\end{leansources}

\leanentry{nonzero}{Nonzero actual cup powers}
\begin{leansources}
\item \href{../../../lean/Aoki/Cech/CochainValues.lean}{\nolinkurl{Aoki/Cech/CochainValues.lean}}, line 140:\\*
\nolinkurl{aleph_powerCochain_not_boundary}.
\item \href{../../../lean/Aoki/Cech/NonzeroClass.lean}{\nolinkurl{Aoki/Cech/NonzeroClass.lean}}, line 62:\\*
\nolinkurl{aleph_topPowerClass_ne_zero}.
\item \href{../../../lean/Aoki/Cech/CupNonvanishing.lean}{\nolinkurl{Aoki/Cech/CupNonvanishing.lean}}, line 58:\\*
\nolinkurl{aleph_cupPower_ne_zero}.
\item \href{../../../lean/Aoki/Cech/CupNonvanishing.lean}{\nolinkurl{Aoki/Cech/CupNonvanishing.lean}}, line 66:\\*
\nolinkurl{cupPower_next_eq_zero}.
\end{leansources}

\leanentry{transfinite}{Closure under transfinite extensions}
\begin{leansources}
\item \href{../../../lean/Aoki/Homological/TransfiniteDimension.lean}{\nolinkurl{Aoki/Homological/TransfiniteDimension.lean}}, line 24:\\*
\nolinkurl{extOne_vanish_iff_surjective_comp}.
\item \href{../../../lean/Aoki/Homological/TransfiniteDimension.lean}{\nolinkurl{Aoki/Homological/TransfiniteDimension.lean}}, line 86:\\*
\nolinkurl{extOne_vanish_of_transfinite}.
\item \href{../../../lean/Aoki/Homological/TransfiniteDimension.lean}{\nolinkurl{Aoki/Homological/TransfiniteDimension.lean}}, line 107:\\*
\nolinkurl{ext_succ_vanish_of_transfinite}.
\item \href{../../../lean/Aoki/Homological/TransfiniteDimension.lean}{\nolinkurl{Aoki/Homological/TransfiniteDimension.lean}}, line 134:\\*
\nolinkurl{hasProjectiveDimensionLE_of_transfinite}.
\end{leansources}

\leanentry{cover-bound}{Compact-open covers and projective dimension}
\begin{leansources}
\item \href{../../../lean/Aoki/Sheaf/ProjectiveOpen.lean}{\nolinkurl{Aoki/Sheaf/ProjectiveOpen.lean}}, line 52:\\*
\nolinkurl{projective_freeOpen_countable_iSup}.
\item \href{../../../lean/Aoki/Sheaf/FreeOpenColimit.lean}{\nolinkurl{Aoki/Sheaf/FreeOpenColimit.lean}}, line 57:\\*
\nolinkurl{freeOpenUnionCoconeIsColimit}.
\item \href{../../../lean/Aoki/Sheaf/FreeOpenFiltration.lean}{\nolinkurl{Aoki/Sheaf/FreeOpenFiltration.lean}}, line 74:\\*
\nolinkurl{initialFreeOpenFiltration_continuous}.
\item \href{../../../lean/Aoki/Sheaf/FreeOpenExact.lean}{\nolinkurl{Aoki/Sheaf/FreeOpenExact.lean}}, line 64:\\*
\nolinkurl{freeOpenMayerVietoris_shortExact}.
\item \href{../../../lean/Aoki/Sheaf/FreeOpenExact.lean}{\nolinkurl{Aoki/Sheaf/FreeOpenExact.lean}}, line 70:\\*
\nolinkurl{freeOpenUnionQuotientIso}.
\item \href{../../../lean/Aoki/Sheaf/CardinalDimension.lean}{\nolinkurl{Aoki/Sheaf/CardinalDimension.lean}}, line 54:\\*
\nolinkurl{freeOpen_iSup_projectiveDimensionLE}.
\item \href{../../../lean/Aoki/Sheaf/CardinalDimension.lean}{\nolinkurl{Aoki/Sheaf/CardinalDimension.lean}}, line 111:\\*
\nolinkurl{sheaf_cohomology_eq_zero_of_compactOpen_cover_cardinal}.
\end{leansources}

\leanentry{minimality}{Universal bounds from cardinality and weight}
\begin{leansources}
\item \href{../../../lean/Aoki/Topology/Weight.lean}{\nolinkurl{Aoki/Topology/Weight.lean}}, line 93:\\*
\nolinkurl{exists_point_indexed_compactOpen_cover}.
\item \href{../../../lean/Aoki/Topology/Weight.lean}{\nolinkurl{Aoki/Topology/Weight.lean}}, line 128:\\*
\nolinkurl{mk_compactOpens_le_weight}.
\item \href{../../../lean/Aoki/Sheaf/CardinalDimension.lean}{\nolinkurl{Aoki/Sheaf/CardinalDimension.lean}}, line 149:\\*
\nolinkurl{sheaf_cohomology_eq_zero_of_cardinal_lt_aleph}.
\item \href{../../../lean/Aoki/Sheaf/CardinalDimension.lean}{\nolinkurl{Aoki/Sheaf/CardinalDimension.lean}}, line 161:\\*
\nolinkurl{sheaf_cohomology_eq_zero_of_weight_lt_aleph}.
\end{leansources}

\begin{thebibliography}{9}
\bibitem{Aoki}
Ko Aoki, \emph{On cohomology of locally profinite sets},
\href{https://arxiv.org/abs/2411.05995v1}{arXiv:2411.05995v1} (2024).
All theorem and question numbers quoted here refer to this version.
\bibitem{Mathlib}
The mathlib community, \emph{Mathlib}, version v4.31.0,
\href{https://github.com/leanprover-community/mathlib4/tree/fabf563a7c95a166b8d7b6efca11c8b4dc9d911f}{source at the pinned commit}.
\end{thebibliography}
\end{document}
